\documentclass[11pt,a4paper]{article}

% ---- Packages ----
\usepackage[utf8]{inputenc}
\usepackage[T1]{fontenc}
\usepackage{amsmath,amssymb,amsthm}
\usepackage{mathtools}
\usepackage[margin=1in,headheight=14pt]{geometry}
\usepackage{hyperref}
\usepackage{enumitem}
\usepackage{xcolor}
\usepackage{tcolorbox}
\usepackage{fancyhdr}
\usepackage{booktabs}

% ---- Hyperref ----
\hypersetup{
    colorlinks=true,
    linkcolor=red,
    citecolor=blue,
    urlcolor=blue
}

% ---- Header ----
\pagestyle{fancy}
\fancyhf{}
\fancyhead[L]{\textit{Lesson 6: Metric and Normed Spaces}}
\fancyhead[R]{\thepage}
\renewcommand{\headrulewidth}{0.4pt}

% ---- Theorem Environments ----
\theoremstyle{definition}
\newtheorem{definition}{Definition}[section]
\newtheorem{example}{Example}[section]
\newtheorem{exercise}{Exercise}[section]

\theoremstyle{plain}
\newtheorem{theorem}{Theorem}[section]
\newtheorem{lemma}[theorem]{Lemma}
\newtheorem{proposition}[theorem]{Proposition}
\newtheorem{corollary}[theorem]{Corollary}

\theoremstyle{remark}
\newtheorem{remark}{Remark}[section]

% ---- Colored Boxes ----
\tcbuselibrary{theorems,skins,breakable}

\newtcolorbox{keyresult}[1][]{
    colback=green!5!white,
    colframe=green!50!black,
    fonttitle=\bfseries,
    title=#1,
    breakable
}

\newtcolorbox{rlconnection}[1][]{
    colback=blue!5!white,
    colframe=blue!50!black,
    fonttitle=\bfseries,
    title=#1,
    breakable
}

\newtcolorbox{warning}[1][]{
    colback=red!5!white,
    colframe=red!50!black,
    fonttitle=\bfseries,
    title=#1,
    breakable
}

% ---- Operators ----
\newcommand{\R}{\mathbb{R}}
\newcommand{\N}{\mathbb{N}}
\newcommand{\Q}{\mathbb{Q}}
\newcommand{\Z}{\mathbb{Z}}
\newcommand{\C}{\mathbb{C}}
\newcommand{\Prob}{\mathbb{P}}
\newcommand{\E}{\mathbb{E}}
\newcommand{\indicator}{\mathbf{1}}
\DeclareMathOperator{\diam}{diam}
\DeclareMathOperator{\Int}{int}
\DeclareMathOperator{\cl}{cl}
\DeclareMathOperator{\Span}{span}
\DeclareMathOperator{\id}{id}

% ---- Title ----
\title{Lesson 17: Bounded Linear Operators and Dual Spaces}
\author{Curriculum: A Rigorous Learning Path to Multi-Agent Deep RL}
\date{}

\begin{document}
\maketitle
\tableofcontents
\newpage

\setcounter{section}{4}

\section{Bounded Linear Operators}
\label{sec:operators}

\subsection{Definitions and Basic Properties}

\begin{definition}[Linear Operator]
\label{def:linear_op}
Let $X$ and $Y$ be normed spaces. A map $T \colon X \to Y$ is a \emph{linear
operator} if $T(\alpha x + \beta y) = \alpha T(x) + \beta T(y)$ for all
$x, y \in X$ and scalars $\alpha, \beta$.
\end{definition}

\begin{definition}[Bounded Linear Operator]
\label{def:bounded_op}
A linear operator $T \colon X \to Y$ is \emph{bounded} if there exists $c \ge 0$
such that $\|Tx\|_Y \le c \|x\|_X$ for all $x \in X$. The \emph{operator norm} is
\[
\|T\| = \sup_{x \ne 0} \frac{\|Tx\|_Y}{\|x\|_X} = \sup_{\|x\|_X = 1} \|Tx\|_Y
= \sup_{\|x\|_X \le 1} \|Tx\|_Y.
\]
\end{definition}

\begin{proposition}[Equivalent Formulations of the Operator Norm]
\label{prop:op_norm_equiv}
The three suprema in Definition~\ref{def:bounded_op} are equal.
\end{proposition}

\begin{proof}
Let $\alpha = \sup_{x \ne 0} \|Tx\|/\|x\|$, $\beta = \sup_{\|x\|=1} \|Tx\|$,
$\gamma = \sup_{\|x\| \le 1} \|Tx\|$.

$\alpha = \beta$: For $x \ne 0$, set $u = x/\|x\|$, so $\|u\| = 1$ and
$\|Tx\|/\|x\| = \|Tu\|$. Thus $\sup_{x \ne 0} \|Tx\|/\|x\| =
\sup_{\|u\|=1} \|Tu\|$.

$\beta \le \gamma$: The set $\{\|x\| = 1\}$ is contained in $\{\|x\| \le 1\}$.

$\gamma \le \alpha$: For $\|x\| \le 1$ with $x \ne 0$,
$\|Tx\| = \|x\| \cdot \|Tx\|/\|x\| \le 1 \cdot \alpha = \alpha$. Also $\|T(0)\| = 0 \le \alpha$.
So $\gamma \le \alpha = \beta \le \gamma$, giving equality.
\end{proof}

\begin{theorem}[Bounded $\Leftrightarrow$ Continuous]
\label{thm:bounded_iff_cts}
A linear operator $T \colon X \to Y$ between normed spaces is bounded if and only
if it is continuous.
\end{theorem}

\begin{proof}
($\Rightarrow$) If $T$ is bounded, then $\|Tx - Ty\| = \|T(x-y)\| \le \|T\| \|x - y\|$,
so $T$ is Lipschitz, hence continuous.

($\Leftarrow$) Suppose $T$ is continuous but not bounded. Then for each $n$, there
exists $x_n$ with $\|x_n\| = 1$ and $\|Tx_n\| > n$. Set $y_n = x_n/n$. Then
$\|y_n\| = 1/n \to 0$, so $y_n \to 0$, but $\|Ty_n\| = \|Tx_n\|/n > 1$. This
contradicts the continuity of $T$ at $0$ (since $y_n \to 0$ should imply
$Ty_n \to T(0) = 0$).
\end{proof}

\begin{remark}
For linear operators, continuity at a single point (say $0$) implies continuity
everywhere, because $T(x_n) - T(x) = T(x_n - x)$ and $x_n - x \to 0$.
\end{remark}

\begin{warning}[Unbounded Operators Exist]
Not every linear operator is bounded. On the space of polynomials $P[0,1]$ (with
the sup norm), the differentiation operator $D(f) = f'$ is linear but unbounded:
$\|x^n\|_\infty = 1$ but $\|D(x^n)\|_\infty = \|n x^{n-1}\|_\infty = n \to \infty$.
\end{warning}

\subsection{The Space $B(X,Y)$}

\begin{definition}[Space of Bounded Operators]
\label{def:BXY}
For normed spaces $X, Y$, we write $B(X,Y)$ for the set of all bounded linear
operators $T \colon X \to Y$, equipped with the operator norm. When $Y = X$, we
write $B(X) = B(X,X)$.
\end{definition}

\begin{theorem}[$B(X,Y)$ is a Normed Space]
\label{thm:BXY_normed}
$B(X,Y)$ is a normed space under the operator norm.
\end{theorem}

\begin{proof}
$B(X,Y)$ is a vector space: for $T, S \in B(X,Y)$ and scalar $\alpha$,
$(T + S)(x) = Tx + Sx$ and $(\alpha T)(x) = \alpha Tx$ are bounded linear operators
(with $\|T + S\| \le \|T\| + \|S\|$ and $\|\alpha T\| = |\alpha| \|T\|$).

The operator norm satisfies the norm axioms:
\begin{itemize}[nosep]
\item \ref{N1}: $\|T\| = 0 \iff \|Tx\| = 0$ for all $\|x\| \le 1 \iff Tx = 0$
for all $x \iff T = 0$.
\item \ref{N2}: $\|\alpha T\| = \sup_{\|x\| \le 1} \|\alpha Tx\| =
|\alpha| \sup_{\|x\| \le 1} \|Tx\| = |\alpha| \|T\|$.
\item \ref{N3}: $\|T + S\| = \sup_{\|x\| \le 1} \|(T+S)x\| \le
\sup_{\|x\| \le 1} (\|Tx\| + \|Sx\|) \le \|T\| + \|S\|$.
\end{itemize}
\end{proof}

\begin{theorem}[$B(X,Y)$ is Banach When $Y$ is Banach]
\label{thm:BXY_banach}
If $Y$ is a Banach space, then $B(X,Y)$ is a Banach space.
\end{theorem}

\begin{proof}
Let $(T_n)$ be a Cauchy sequence in $B(X,Y)$. For each $x \in X$,
$\|T_n x - T_m x\| \le \|T_n - T_m\| \|x\| \to 0$, so $(T_n x)$ is Cauchy in $Y$.
Since $Y$ is Banach, the limit exists: define $Tx = \lim_{n \to \infty} T_n x$.

\textbf{$T$ is linear.} $T(\alpha x + \beta y) = \lim T_n(\alpha x + \beta y) =
\lim (\alpha T_n x + \beta T_n y) = \alpha Tx + \beta Ty$.

\textbf{$T$ is bounded.} Since $(T_n)$ is Cauchy, it is bounded: $\|T_n\| \le M$
for some $M$. For each $x$, $\|Tx\| = \lim \|T_n x\| \le M \|x\|$, so
$\|T\| \le M$.

\textbf{$T_n \to T$ in operator norm.} Given $\varepsilon > 0$, choose $N$ such
that $\|T_n - T_m\| < \varepsilon$ for $n, m \ge N$. For any $x$ with $\|x\| \le 1$
and $n \ge N$:
$\|T_n x - T_m x\| \le \varepsilon$ for all $m \ge N$. Letting $m \to \infty$:
$\|T_n x - Tx\| \le \varepsilon$. Taking the sup over $\|x\| \le 1$:
$\|T_n - T\| \le \varepsilon$ for $n \ge N$.
\end{proof}

\subsection{Important Examples}

\begin{example}[Matrix Operators]
\label{ex:matrix_op}
Every $m \times n$ real matrix $A$ defines a bounded linear operator
$T_A \colon (\R^n, \|\cdot\|_p) \to (\R^m, \|\cdot\|_p)$ by $T_A(x) = Ax$.
Boundedness follows from finite-dimensionality (all linear operators on
finite-dimensional spaces are bounded, as we prove below).
\end{example}

\begin{proposition}[All Linear Operators on Finite-Dimensional Spaces are Bounded]
\label{prop:finite_dim_bounded}
If $X$ is a finite-dimensional normed space and $Y$ is any normed space, then
every linear operator $T \colon X \to Y$ is bounded.
\end{proposition}

\begin{proof}
Let $\{e_1, \ldots, e_n\}$ be a basis for $X$. For $x = \sum_{i=1}^n \alpha_i e_i$,
by the equivalence of norms (Theorem~\ref{thm:norm_equiv}), there exists $c > 0$
with $\sum |\alpha_i| \le c \|x\|_X$ (since $x \mapsto \sum |\alpha_i|$ is a norm
on $X$, equivalent to $\|\cdot\|_X$). Then:
\[
\|Tx\| = \Bigl\|\sum_{i=1}^n \alpha_i Te_i\Bigr\| \le \sum_{i=1}^n |\alpha_i| \|Te_i\|
\le c \|x\|_X \cdot \max_i \|Te_i\|.
\]
So $T$ is bounded with $\|T\| \le c \cdot \max_i \|Te_i\|$.
\end{proof}

\begin{example}[Integral Operator]
\label{ex:integral_op}
Let $k \in C([a,b] \times [a,b])$ and define $T \colon C[a,b] \to C[a,b]$ by
\[
(Tf)(s) = \int_a^b k(s,t)\, f(t)\, dt.
\]
This is a bounded linear operator with
$\|T\| \le \max_{s \in [a,b]} \int_a^b |k(s,t)|\, dt$.
Indeed, $|(Tf)(s)| \le \int_a^b |k(s,t)| |f(t)|\, dt \le \|f\|_\infty \int_a^b |k(s,t)|\, dt$.
\end{example}

\begin{example}[Right Shift Operator]
\label{ex:shift}
On $\ell^p$ ($1 \le p \le \infty$), the right shift $R(x_1, x_2, \ldots) =
(0, x_1, x_2, \ldots)$ is a bounded linear operator with $\|R\| = 1$: clearly
$\|Rx\|_p = \|x\|_p$ for all $x$.
\end{example}

\begin{rlconnection}[The Bellman Operator is Bounded]
The Bellman expectation operator $T^\pi \colon b(\mathcal{S}) \to b(\mathcal{S})$
defined by $(T^\pi V)(s) = r^\pi(s) + \gamma \sum_{s'} P^\pi(s'|s) V(s')$ is a
bounded (in fact, affine) operator on the Banach space of bounded functions. Its
``linear part'' $V \mapsto \gamma \sum_{s'} P^\pi(s'|s) V(s')$ has operator norm
at most $\gamma < 1$ (under $\|\cdot\|_\infty$), which is why $T^\pi$ is a
contraction. The framework of $B(X,Y)$ and operator norms gives us the precise
language to state and prove this.
\end{rlconnection}

\subsection{Dual Spaces}

\begin{definition}[Dual Space]
\label{def:dual}
The \emph{(continuous) dual space} of a normed space $X$ is
$X^* = B(X, \R)$ (or $B(X, \C)$ if $X$ is a complex space): the space of all
bounded linear functionals on $X$, equipped with the operator norm
$\|f\|_{X^*} = \sup_{\|x\| \le 1} |f(x)|$.
\end{definition}

\begin{corollary}
\label{cor:dual_banach}
For any normed space $X$, the dual space $X^*$ is a Banach space.
\end{corollary}

\begin{proof}
$\R$ (or $\C$) is a Banach space, so $X^* = B(X, \R)$ is Banach by
Theorem~\ref{thm:BXY_banach}.
\end{proof}

\begin{remark}
The dual space is always complete, even if $X$ itself is not.
\end{remark}

\begin{example}[Dual of $\ell^p$]
\label{ex:dual_lp}
For $1 \le p < \infty$ with $1/p + 1/q = 1$ (where $q = \infty$ when $p = 1$),
$(\ell^p)^* \cong \ell^q$ isometrically: every bounded linear functional on
$\ell^p$ has the form $f(x) = \sum_{n=1}^\infty y_n x_n$ for a unique
$y = (y_n) \in \ell^q$, with $\|f\| = \|y\|_q$. This is the Riesz representation
theorem for $\ell^p$ spaces.
\end{example}

%% ============================================================
%% SECTION 6: COMPACTNESS
%% ============================================================


\end{document}